% DL9MJ, Matthias Jung, 2026
\begin{tikzpicture}

    % Parameter
    \pgfmathsetmacro{\f}{1.4}      % Brennweite
    \pgfmathsetmacro{\ap}{2.4}     % halbe Apertur
    \pgfmathsetmacro{\xend}{5.0}   % Ende der reflektierten Strahlen

    % Brennpunkt
    \coordinate (F) at (\f,0);

    % Parabolspiegel: y^2 = 4 f x
    \draw[ultra thick]
        plot[
            variable=\t,
            domain=-\ap:\ap,
            samples=100,
            smooth
        ]
        ({\t*\t/(4*\f)},{\t});
    
    \draw({(-\ap)*(-\ap)/(4*\f)},{-\ap}) coordinate(end);

    % Feed / Antenne im Brennpunkt
    \node[bareantenna, scale=1.0, anchor=west] (ant) at (F) {};

    \draw(end) to [open] ++(2,0) node[trxshape, scale=1, anchor=west](trx){};

    % Anschlussleitung und TRX
    \draw (ant)
        -- ++(0,-1.2) coordinate(h)
        -| (trx.north);

    % Strahlen vom Brennpunkt zum Spiegel und danach parallel
    \foreach \yy in {-2.0,-1.2,-0.4,0.4,1.2,2.0} {
        \pgfmathsetmacro{\xx}{\yy*\yy/(4*\f)}

        % vom Brennpunkt zum Spiegel
        \draw[DARCred, thick]
            (F) -- (\xx,\yy);

        % reflektierter Strahl parallel zur Achse
        \draw[DARCred, thick]
            (\xx,\yy) -- (\xend,\yy);
    }
\end{tikzpicture}